\documentclass[tikz]{standalone}
\usepackage{bm}
\usepackage{amsmath}
\usepackage{amssymb}
\usetikzlibrary{calc}

\begin{document}

\usepgflibrary{shadings} 

\pgfdeclareradialshading{myflameshade}{\pgfpoint{0}{0.05cm}}{%
  color(0.0cm)=(white); color(0.5cm)=(red!40!yellow);
  color(0.7cm)=(red!50!yellow); color(0.9cm)=(white)
}

\tikzset{
  vial/.cd,
  water height/.initial=0.74,
  water color/.initial=cyan!60!blue!20!white,
  water color bottom/.initial=cyan!60!blue!20!white!60!black,
  %
  sand peaker/.is choice,
  sand peaker/true/.code={\def\issandpeaker{1}},
  sand peaker/false/.code={\def\issandpeaker{0}},
  sand peaker=true,
  %
  /tikz/flame/.pic={
    \def\wflame{1.0}\def\hflame{6.0}
    % transform shape here ensures the flame itself can be scaled
    \begin{scope}[transform shape]
      \clip (0,-0.1*\hflame) to[out=0, in=-90] (\wflame,0) 
            arc (0:180:\wflame\space and \hflame) to[out=-90,in=180] cycle;
      \shade[shading=myflameshade] (0,0) ellipse ({\wflame} and {\hflame});
    \end{scope}
  },
  %
  /tikz/vial/.pic={
    \tikzset{vial/.cd, #1}
    \pgfkeysgetvalue{/tikz/vial/water height}{\ratio}
    \pgfkeysgetvalue{/tikz/vial/water color}{\vcol}
    \pgfkeysgetvalue{/tikz/vial/water color bottom}{\vcolbot}
    
    % Geometry Constants
    \def\Rx{1.1} \def\Ry{0.2} \def\H{2.6}
    \def\h{\ratio*\H}
    \def\pt{0.07*\H} \def\prx{0.040*\H} \def\pry{0.006*\H} \def\pl{0.25*\H}
    
    % transform shape is the magic key here
    \begin{scope}[transform shape]
        \colorlet{mydarkblue}{blue!50!black}
        \colorlet{darkwatercol}{\vcol}
        \colorlet{watercolorbottom}{\vcolbot}
        \colorlet{metalcol}{blue!40!black!10!white}
        \tikzstyle{dark water}=[draw=mydarkblue, top color=darkwatercol, bottom color=watercolorbottom]
        \tikzstyle{metal}=[draw=metalcol!20!black, top color=metalcol, bottom color=metalcol!90!black]

        
        \draw[dark water] (-\Rx,\h) -- (-\Rx,0) arc (180:360:{\Rx} and {\Ry}) -- (\Rx,\h);
        \draw[metal] (-\Rx,\h) --++ (0,-\pt) arc (180:360:{\Rx} and {\Ry}) --++ (0,\pt);
        \draw[metal] (0,\h) ellipse ({\Rx} and {\Ry});
        \draw[metal] (-\prx,\h+\pl) --++ (0,-\pl) arc (180:360:{\prx} and {\pry}) --++ (0,\pl);
        \draw[metal](0,\h+\pl) ellipse ({\prx} and {\pry});

        % sand  
        \fill[orange!40!white] ({-0.6*(1-0.5*\ratio)},\h) .. controls (-0.0,\h+0.3*\H-0.3*\h+0.4) and (0.3,\h+0.3*\H-0.3*\h+0.4) .. ({0.6*(1-0.5*\ratio)},\h) coordinate[pos=0.55] (A) to[bend left=20] cycle;

        % flowing sand and sand peaker
        \ifnum\issandpeaker=1
        \draw[ultra thick, orange!40!white] (A.center) to [bend left=10](0.8,4);
        \fill[orange!40!white] (0.8,4)--(2,3)--(2.4,3.4)--(1.4,4.3)--cycle;
        \draw[rounded corners, very thick] (0.8,4)--(2,3)--(2.4,3.4)--(1.2,4.5);
        \fi
        % top edge
        \draw[thick] (-\Rx,\H) -- (-\Rx,0) arc (180:360:{\Rx} and {\Ry}) -- (\Rx,\H) (0,\H) ellipse ({\Rx} and {\Ry});

    \end{scope}
  }
}


% a new figure of heating against a constant force
\begin{tikzpicture}

    \def\scale{1.3};
    \def\H{2.6};
    \def\dH{0.25*\H};

  % initial time
  %\pic[scale=0.3,anchor=center] at (0, -1.5) {flame};
  \pic[scale=\scale] at (0,0) {vial={water height=0.8, sand peaker=false}};
  %\node[anchor=center] at (0,0.55) {$(p_0,T_0,v_0,u_0)$};
  %\node[anchor=center] at (0,0.2) {$\bm{v}=\bm{0}$};
  \node[anchor=center] at (0,0.2) {Initial};

  \node[anchor=center] at (2,1) {$\dots$};
  
  % heating
  \pic[scale=0.3,anchor=center] at (4, -1.5) {flame};
  \pic[scale=\scale] at (4,0) {vial={water height=0.6, water color=orange!20!white, water color bottom=orange!20!white}};
  %\node[anchor=center] at (4,0.55) {$(p_i,T_i,v_i,u_i)$};
  %\node[anchor=center] at (4,0.2) {$\bm{v}=\bm{0}$};

  \node[anchor=center] at (6,1) {$\dots$};
  
  % final time
  %\pic[scale=0.3,anchor=center] at (8, -1.5) {flame};
  \pic[scale=1.3] at (8,0) {vial={water height=0.4, water color=red!80!orange!40!white, water color bottom=red!80!orange!40!white, sand peaker = false}};
  %\node[anchor=center] at (8,0.55) {$(p_N,T_N,v_N,u_N)$};
  %\node[anchor=center] at (8,0.2) {$\bm{v}=\bm{0}$};
  \node[anchor=center] at (8,0.2) {Final};
   
\end{tikzpicture}

\end{document}
